%% =========================================================================
%%  Приложение «Программная реализация и воспроизводимость»
%%  Перенесено из sections/reproducibility.tex 2026-05-01.
%%  В блоке \appendix главы автоматически нумеруются \Asbuk (А, Б, В),
%%  поэтому здесь \chapter (без звёздочки) — заголовок «ПРИЛОЖЕНИЕ X»
%%  формируется переопределённым \@makechapterhead в main.tex.
%% =========================================================================

\chapter{Программная реализация и воспроизводимость}
\label{ch:reproducibility}

\section{Структура репозитория}

Все эксперименты и исходные тексты диссертации собраны в публичном
репозитории \texttt{SP-Broyden}.\footnote{Адрес репозитория и команда
сборки указаны в файле \texttt{README.md} корня проекта; для версии,
зафиксированной на момент защиты, использован тег \texttt{master-thesis-2026}.}
Структура каталога:
\begin{description}\itemsep0pt
  \item[\texttt{mipt\_thesis\_master/}] \LaTeX-источник диссертации
        (\texttt{main.tex}, \texttt{sections/}, \texttt{appendix/}),
        скрипт сборки \texttt{build.sh} (двойной прогон \texttt{pdflatex})
        и статический проверочный скрипт \texttt{check\_static.py},
        сверяющий \verb|\cite|/\verb|\bibitem|, \verb|\ref|/\verb|\label|.
  \item[\texttt{tezisy.ipynb}] эксперименты главы~2 (SP-Broyden на
        нелинейных системах: тест Broyden, проверка ограниченной
        деградации $\Fro{B_k - J(x_k)}$).
  \item[\texttt{basin.ipynb}] эксперименты главы~3 (SS-SR1: бассейн
        сходимости, скейлинг радиуса по $(\alpha,\beta)$).
  \item[\texttt{pics.ipynb}] двумерные иллюстрации (Curved Valley,
        проекция $\Pi_k$, разложение ошибки якобиана).
  \item[\texttt{run\_seeds.py}] мульти-seed эксперимент главы~4
        (SP-AFD vs SP-Broyden vs VIQA-Broyden, рисунок
        \ref{fig:viji_seeds}); сохраняет промежуточные траектории
        в \texttt{results.npz}.
\end{description}

\section{Окружение}

Эксперименты прогонялись в окружении:
\begin{itemize}\itemsep0pt
  \item Python~3.10 (CPython, x86\_64 Linux);
  \item NumPy~1.26 (линейная алгебра, генератор
        \texttt{numpy.random.default\_rng}),
        SciPy~1.11 (только для проверки референсного решения),
        Matplotlib~3.8 (графики);
  \item \LaTeX: TeX~Live~2025+ с поддержкой кириллицы (\texttt{babel-russian},
        \texttt{T2A}-энкодинг, пакет \texttt{mathtext} класса
        \texttt{mipt-thesis-bs}).
\end{itemize}
Для воспроизведения PDF достаточно команды (двойной \texttt{pdflatex}):
\begin{center}
\verb|cd mipt_thesis_master && ./build.sh|
\end{center}

\section{Случайные числа и seed'ы}

Все случайные величины генерируются исключительно через современный
интерфейс \texttt{rng = numpy.random.default\_rng(seed)} (PCG64) с
явно передаваемым целочисленным seed'ом~--- глобальный
\texttt{numpy.random.seed} не используется, чтобы исключить скрытые
зависимости между прогонами.
\begin{itemize}\itemsep0pt
  \item Генерация задачи (главы~2 и~4):
        матрица $U$~--- из QR-разложения $\mathcal N(0,1)^{n\times n}$,
        сдвиг $c\sim\mathcal N(0,0{,}3^2)$, целевой корень $x^{*}$~---
        $\mathcal N(0,0{,}5^2)$. Дальнейшая нормировка $c$ из условия
        $F(x^{*})=0$ приведена в \texttt{run\_seeds.py}, функция
        \texttt{make\_problem}.
  \item Стартовая точка (главы~2 и~4): $x_0 = x^{*} + \sigma\cdot \mathcal N(0,1)^n$,
        $\sigma=0{,}3$ для $n=30$ (даёт $\Norm{x_0-x^{*}}\approx 1{,}6$).
  \item Эксперимент с доверительными интервалами (рис.~\ref{fig:viji_seeds}):
        seed'ы постановки задачи $\{0,1,\ldots,9\}$, seed'ы стартовой
        точки $\{1000,\ldots,1009\}$. Полный набор фиксирован в
        \texttt{run\_seeds.py} (\texttt{SEEDS = list(range(10))}).
  \item Эксперименты главы~3 (\texttt{basin.ipynb}): seed
        \texttt{0}~--- сгенерирован один раз для каждой пары
        $(\alpha,\beta)\in\{10,50,100\}\times\{1,1{,}5,2\}$;
        вариативность по seed'ам не агрегировалась
        (см.\ ограничения в \S\ref{sec:viji:discussion}).
  \item Статистические эксперименты главы~3
        (рис.~\ref{fig:ndim_stat_sr1}, \ref{fig:ndim_stat_psb},
        \ref{fig:ndim_stat_rpast_psb}):
        seed \texttt{20\,260\,503} (PCG64), $50$~единичных направлений
        на сфере $S^{n-1}$ для каждой размерности $n\in\{10,20,50\}$,
        paired-design (одни и те же направления для всех методов
        внутри одной пары SS-X / X). Из~каждой траектории сохраняется
        как $\Norm{\nabla f(x_k)}_2$ (для convergence-фигур), так и
        $\Norm{B_k S_p - Y_p}_F$ (rpast-метрика для~PSB-пары~--- прямая
        иллюстрация теоремы~\ref{thm:ss_props}, п.~3; для~SR1-пары
        rpast-данные сохраняются в~\texttt{.npz} без построения фигуры
        в~главе, см.~\S\ref{par:ssn-rpast}).
\end{itemize}

\section{Единый критерий остановки}

Поскольку три класса задач~--- нелинейные системы, безусловная
оптимизация и сильно монотонные VI~--- допускают разные естественные
индикаторы сходимости, мы документируем критерии явно. Везде
численный допуск $\varepsilon_{\mathrm{tol}}=10^{-10}$, максимальное
число итераций фиксируется в зависимости от метода.
\begin{itemize}\itemsep0pt
  \item \textbf{Глава~2 (нелинейные системы $F(x)=0$).}
        Останов по невязке $\Norm{F(x_k)}_2 \le \varepsilon_{\mathrm{tol}}$.
        Этот критерий совпадает с критерием класса \texttt{tezisy.ipynb}
        (строка \verb|tol = 1e-10|).
  \item \textbf{Глава~3 (минимизация $\min f(x)$).}
        Останов по норме градиента $\Norm{\nabla f(x_k)}_2 \le \varepsilon_{\mathrm{tol}}$
        в экспериментах сходимости; для построения бассейна сходимости
        в \texttt{basin.ipynb} используется радиальный критерий
        $\Norm{x_k - x^{*}}_2 \le \delta_{\mathrm{basin}}$ (с
        $\delta_{\mathrm{basin}} = 10^{-6}$ и фиксированным бюджетом
        итераций), поскольку искомая величина~--- именно радиус,
        с которого метод сошёлся.
  \item \textbf{Глава~4 (сильно монотонная VI с оператором $F$).}
        В качестве индикатора сходимости используется радиус
        $\Norm{x_k - x^{*}}_2$ (эталонный корень $x^{*}$ известен из
        конструкции тестовой задачи: $c$ выбрана так, чтобы $F(x^{*})=0$).
        Это эквивалентно $\Norm{F(x_k)}_2$ с точностью до константы:
        из сильной монотонности $\mu\Norm{x_k-x^{*}}_2\le\Norm{F(x_k)}_2\le L_1\Norm{x_k-x^{*}}_2$,
        откуда оба индикатора пропорциональны радиусу с точностью
        до $\mathcal O(L_1/\mu)$. На прикладных задачах,
        где $x^{*}$ не задан явно, в качестве индикатора рекомендуется
        стандартный gap-функционал
        $\Phi_R(x)=\sup_{\Norm{y-x}\le R}\langle F(y),x-y\rangle$
        (см., например, \cite[\S 5]{Agafonov2024VIJI}).
\end{itemize}

\section{Соответствие между текстом и ноутбуками}

\begin{center}
\begin{tabularx}{\textwidth}{@{}XX@{}}
\toprule
Раздел текста & Источник эксперимента \\
\midrule
\S\ref{sec:viji:numerics}, рис.~\ref{fig:viji_conv} (одиночный прогон) & \texttt{tezisy.ipynb}, ячейка <<SP-AFD vs SP-Broyden vs VIQA>> \\
\S\ref{sec:viji:numerics}, рис.~\ref{fig:viji_seeds} (10 seed'ов, IQR) & \texttt{run\_seeds.py} \\
Глава~2, проверка $\Fro{E_k}=\Fro{B_k - J(x_k)}$ & \texttt{tezisy.ipynb}, ячейка <<bounded deterioration>> \\
Глава~2, 2D-иллюстрации проектора $\Pi_k$ & \texttt{pics.ipynb} \\
\S\ref{sec:ss_numerics}, рис.~\ref{fig:ndim_stat_sr1}, \ref{fig:ndim_stat_psb}, \ref{fig:ndim_stat_rpast_psb} (50~стартов, $n\in\{10,20,50\}$; convergence + $\Norm{B_k S_p - Y_p}_F$ для PSB-пары) & \texttt{diag\_ndim\_stat.py} \\
\bottomrule
\end{tabularx}
\end{center}
